
\chapter[ML classifiers: Na\"{i}ve Bayes (3 of 3)]
{\huge ML classifiers: Na\"{i}ve Bayes (3 of 3)}
\label{ch:naiveBayes3}

\section{Numeric variables}

\index{Jedi}
\index{Sith}

In all the examples above, we had categorical features. Computing something
like ``$P(\textrm{calm}\ |\ \textrm{Jedi})$'' is thus just a matter of counting
the \texttt{Jedi} rows which (exactly) match the value \texttt{calm}. But what
if we have a numeric feature, like a person's age or salary?

One quick-and-dirty way to deal with numeric features is to ``bin'' them, kind
of like a histogram. Note that this is actually what we already did with the
``\texttt{height}'' column in the Jedi example: instead of heights like 4'9"
and 6'6", we had values like \texttt{short} and \texttt{tall}.

This has obvious drawbacks. For one, we're throwing away all the detail
present in the measurement and treating different values exactly alike. For
another, it's up to us to somewhat arbitrarily choose how many bins to use and
where the demarcation lines are. This has an \textit{ad hoc} feel to it.

\subsection{Gaussian Na\"{i}ve Bayes}
\index{Gaussian Na\"{i}ve Bayes}
\index{normal distribution}
\index{Gaussian (normal) distribution}
\index{bell-curvy@``bell-curvy''}
\index{probability density}

Better is to treat the continuous variable as what it really is: a number.
\textbf{Gaussian Na\"{i}ve Bayes} is probably the most common way to deal with
this in a principled fashion: we make the assumption that the numeric feature
is distributed \textbf{normally} (``normal''==``Gaussian''), and estimate its
mean and standard deviation from the training data. Then, to estimate
$P(\textrm{value}\ |\ \textrm{target\_label})$ for any particular test data
point's value, we choose the estimated value of its \textbf{probability
density} (p.~\pageref{probDensity}) and use that.

Remember from Section~\ref{GaussianDistro} (p.~\pageref{GaussianDistro}) that
if a variable is ``distributed normally'' it means that its values are
``bell-curvy.'' There's some average value which occurs a lot, and values
farther and farther from that average are less and less likely to occur.

Now suppose we had a numeric \texttt{height} feature (measured in inches)
instead of a categorical one. Our training data (in the ``\texttt{t}''
\texttt{DataFrame}) might look like this:

\begin{Verbatim}[fontsize=\small,samepage=true,frame=none]
                    face  height  demeanor lightsaber  type
                   clean    62.5  menacing      green  Sith
                   beard    68.3      calm      green  Jedi
                   clean    78.1  menacing        red  Sith
                   beard    63.9      calm       blue  Jedi
                   beard    61.9      calm       blue  Jedi
                   clean    68.2      calm       blue  Sith
                   clean    80.1  menacing        red  Jedi
                   clean    33.5      calm      green  Jedi
                                . . .
\end{Verbatim}

\index{normal distribution}
\index{Gaussian (normal) distribution}

Our Gaussian Na\"{i}ve Bayes assumption will be that \textit{for each class,
the \texttt{height} feature is distributed normally}. Given this assumption, we
can characterize its normal distribution with just two parameters: its
estimated mean and standard deviation. So we take all the \texttt{Jedi} rows,
and calculate their \texttt{height}s' mean and standard deviation:

\begin{Verbatim}[fontsize=\small,samepage=true,frame=single,framesep=3mm]
t[t.type=='Jedi'].height.mean()
t[t.type=='Jedi'].height.std()
\end{Verbatim}
\vspace{-.3in}

\begin{Verbatim}[fontsize=\small,samepage=true,frame=leftline,framesep=5mm,framerule=1mm]
64.0571428571429
9.661931704216865
\end{Verbatim}

\vspace{.1in}
Performing the same calculation for the \texttt{Sith} rows gives us:

\begin{Verbatim}[fontsize=\small,samepage=true,frame=single,framesep=3mm]
t[t.type=='Sith'].height.mean()
t[t.type=='Sith'].height.std()
\end{Verbatim}
\vspace{-.3in}

\begin{Verbatim}[fontsize=\small,samepage=true,frame=leftline,framesep=5mm,framerule=1mm]
68.3902173913043
9.319673312968604
\end{Verbatim}

Interesting. Looks like Sith Lords are a bit taller than Jedi in general, and
slightly less variable.

\vspace{.1in}

\index{frozen distribution@``frozen'' distribution (SciPy)}
\index{spycam}
\index{security camera}

Now let's say we're doing a Jedi-vs.-Sith prediction for a person whose height
(according to our security camera) is 69 inches. We'll need to know the value
of the density at $x$=69 for both our estimated Jedi and estimated Sith
distributions. Using the SciPy library, we can get ``frozen'' distributions for
each:

\index{import@\texttt{import}}

\begin{Verbatim}[fontsize=\small,samepage=true,frame=single,framesep=3mm]
import scipy.stats
jedi_height = scipy.stats.norm(64.0571428571429, 9.661931704216865)
sith_height = scipy.stats.norm(68.3902173913043, 9.319673312968604)
\end{Verbatim}

\index{pdf@\texttt{.pdf()} method}

Think of a frozen distribution as just a variable on which we can call a
\texttt{.pdf()} method to find the height of its distribution at a particular
value. Let's see how high the estimated Jedi and Sith densities are for a
height of 69 inches:

\begin{Verbatim}[fontsize=\small,samepage=true,frame=single,framesep=3mm]
jedi_height.pdf(69)
sith_height.pdf(69)
\end{Verbatim}
\vspace{-.2in}

\begin{Verbatim}[fontsize=\small,samepage=true,frame=leftline,framesep=5mm,framerule=1mm]
0.036230663486274517
0.042717943988373371
\end{Verbatim}

Not surprisingly, the Sith number comes out a fair bit higher, since according
to our estimated density, the average Sith's height is taller than the average
Jedi's.

Keep in mind that the ``.0362'' and ``.0427'' numbers here are a bit hard to
interpret. They are \textit{not} ``the probability that a Jedi/Sith is exactly
69 inches tall'' -- that probability, as with all continuous random variables,
is \textit{zero}. Instead, it's ``the value of the density at $x$=69'' which
isn't really in any units that make sense. It's okay, though, since we are
throwing away out denominator and doing relative comparisons anyway.

\index{normal random value}
\index{random value}
\index{quantiles}
\index{mean}
\index{standard deviation}

By the way, there are also other useful things you can call on a frozen SciPy
distribution, including \texttt{.cdf()} to get the value of the cumulative
density function at any point, \texttt{.ppf()} to get a quantile,
\texttt{.rvs()} to generate random values from it, ordinary statistical
things like \texttt{.median()}, \texttt{.mean()}, \texttt{.std()}, and other
more exotic statistical things like moments, confidence intervals, differential
entropy, survival functions, \textit{etc.} See the documentation for details.

\subsection{Incorporating into the classifier}

Let's add this feature to our classifier. First, the training code for the
not-so-generally-coded version (p.~\pageref{calcConditionals}), which had lines
like this:

\footnotesize
\begin{Verbatim}[fontsize=\small,samepage=true,frame=single,framesep=3mm]
probs['tall|Sith'] = (len(t[(t.type=='Sith') & (t.demeanor=='tall')]) /
    len(t[(t.type=='Sith')]))
\end{Verbatim}
\normalsize

Since \texttt{height} is now numeric, we need to replace all those
height-related lines with code like this:

\begin{Verbatim}[fontsize=\small,samepage=true,frame=single,framesep=3mm]
sith_heights = t[t.type == 'Sith'].height
sith_height = scipy.stats.norm(sith_heights.mean(),sith_heights.std())
jedi_heights = t[t.type == 'Jedi'].height
jedi_height = scipy.stats.norm(jedi_heights.mean(),jedi_heights.std())
\end{Verbatim}

\index{frozen distribution@``frozen'' distribution (SciPy)}
\index{predict@\texttt{predict()}}

``\texttt{sith\_height}'' and ``\texttt{jedi\_height}'' are now our frozen
distribution functions. In \texttt{predict()}, then, we simply use it as
another multiplier:

\begin{Verbatim}[fontsize=\small,samepage=true,frame=single,framesep=3mm]
def predict(face, height, demeanor, lightsaber):
    jediness = (probs[face+"|Jedi"] *
                jedi_height.pdf(height) *               <-- CHANGED
                probs[demeanor+"|Jedi"] *
                probs[lightsaber+"|Jedi"] *
                probs["Jedi"])
    sithness = (probs[face+"|Sith"] *
                sith_height.pdf(height) *               <-- CHANGED
                probs[demeanor+"|Sith"] *
                probs[lightsaber+"|Sith"] *
                probs["Sith"])
    if jediness > sithness:
        return "Jedi"
    else:
        return "Sith"
\end{Verbatim}

Squint at that code and notice the changed lines from
Figure~\ref{fig:predictFunc} (on p.~\pageref{fig:predictFunc}). We've removed
the ``\texttt{probs[height+"|Jedi"]}'' lines, since they're no longer
applicable: \texttt{height} is numeric now. Instead, we multiply our
\texttt{jediness}/\texttt{sithness} by the value of the corresponding class's
estimated density at that point.

\index{dtype@\texttt{dtype}}

To incorporate this numeric variable into the more general version of the code
(p.\pageref{codingNbGenerally}), we need to intelligently distinguish between
numeric and categorical columns. This can be done via the \texttt{.dtype}
accessor on a \texttt{DataFrame} column:

\begin{Verbatim}[fontsize=\small,samepage=true,frame=single,framesep=3mm]
print(t.lightsaber.dtype)
print(t.height.dtype)
\end{Verbatim}
\vspace{-.2in}

\begin{Verbatim}[fontsize=\small,samepage=true,frame=leftline,framesep=5mm,framerule=1mm]
dtype('O')
dtype('float64')
\end{Verbatim}

\index{object@``\texttt{object}'' (a.k.a. \texttt{str}) \texttt{dtype}}
\index{issubdtype@\texttt{issubdtype()}}

(The ``\texttt{O}'' stands for ``object,'' which is Pandas' non-intuitive way
of telling you that the column is something non-numeric, like a string.) You
can directly compare the result to a value like
``\texttt{np.dtype("float64")}'', but the safer and preferred way is to use
NumPy's \texttt{issubdtype()} function which more intelligently detects any
numeric type (\texttt{int}s, \texttt{float}s, \textit{etc.}, no matter the
precision).

Our new generic implementation begins with:

\begin{Verbatim}[fontsize=\footnotesize,samepage=true,frame=single,framesep=3mm]
probs = {}
densities = {}
target = t.columns[len(t.columns)-1]
num_fs= [ c for c in t.columns[:-1] if np.issubdtype(t[c].dtype,np.number) ]
cat_fs = [ c for c in t.columns[:-1] if not np.issubdtype(t[c].dtype,np.number) ]
\end{Verbatim}

We're creating a \texttt{densities} dictionary, somewhat similar to our
\texttt{probs} dictionary, to store feature-specific information for our
numeric variables (\texttt{probs} will continue to store information for
categorical variables). Instead of simple numeric values, \texttt{densities}
will store frozen distribution functions. Additionally, instead of a single
\texttt{features} list like we had on p.~\pageref{codingNbGenerally}, we now
have two lists -- \texttt{num\_fs} and \texttt{cat\_fs} -- for the names of the
numeric and categorical features, respectively.

\index{list!list comprehension}
\index{boxies (square brackets)}
\index{[]@\texttt{[]} (boxies)}

Something else in that code might be new to you: we're using something called a
\textbf{list comprehension}, which is just a sleek way of creating a list
without the need for a \texttt{for} loop. Notice the boxies after the equals
signs for \texttt{num\_fs} and \texttt{cat\_fs}, and the \texttt{if} clauses.
This code says ``if column \texttt{c} is a numeric variable, put it in
\texttt{num\_fs}, otherwise put it in \texttt{cat\_fs}.''

\medskip

All right, now we can populate \texttt{densities} (and \texttt{probs}):

\begin{Verbatim}[fontsize=\small,samepage=true,frame=single,framesep=3mm]
for label in df[target].drop_duplicates():
    rows_with_this_label = df[df[target] == label]
    probs[label] = len(rows_with_this_label)/len(df)
    for feature in numeric_features:
        density = scipy.stats.norm(rows_with_this_label[feature].mean(),
            rows_with_this_label[feature].std())
        densities[feature+"|"+label] = density
    for feature in categorical_features:
        for value in df[feature].drop_duplicates():
            probs[value+"|"+label] = (
                len(df[(df[target] == label) & (df[feature] == value)]) / 
                    len(df[df[target] == label]))
\end{Verbatim}

\index{predict@\texttt{predict()}}
\index{key-value pair}

Our generic \texttt{predict()} function from
(p.\pageref{codingNbGenerallyPredict}) needs to be adjusted somewhat, since if
we just pass a list of values, we can't tell \textit{which} feature a numeric
value is supposed to go with. There are various ways to solve this, but we'll
choose to pass a key/value pair string in this case (like
\texttt{"height=69.3"}). The complete code is in Figure~\ref{generalNbWithNum},
and we call the function like this:

\begin{Verbatim}[fontsize=\small,samepage=true,frame=single,framesep=3mm]
predict(['beard','height=65','menacing','blue'])
predict(['beard','height=85','menacing','blue'])
\end{Verbatim}
\vspace{-.2in}

\begin{Verbatim}[fontsize=\small,samepage=true,frame=leftline,framesep=5mm,framerule=1mm]
['Jedi', 0.85863336769968801]
['Jedi', 0.72742228657164021]
\end{Verbatim}

\begin{figure}[ht]
\centering
\begin{Verbatim}[fontsize=\footnotesize,samepage=true,frame=single,framesep=3mm]
def predict(values):

    scores = []
    top_prob = 0
    top_answer = ""

    for label in df[target].drop_duplicates():

        prob = probs[label]    # Start with the prior

        for value in values:
            if "=" in value:
                # Numeric feature.
                numeric_feature, numeric_value = value.split("=")
                numeric_value = float(numeric_value)
                prob *= densities[numeric_feature+"|"+label].pdf(numeric_value)
            else:
                # Categorical feature.
                prob *= probs[value+"|"+label]

        scores.append(prob)

        if prob > top_prob:
            top_prob = prob
            top_answer = label

    return [ top_answer, top_prob / sum(scores) ]
\end{Verbatim}
\caption{A general Na\"{i}ve Bayes \texttt{predict()} function, enhanced to
allow numeric features (which are passed as key/value strings like
\texttt{"height=51.5"}).}
\label{generalNbWithNum}
\end{figure}

And how does this thing do?

\begin{Verbatim}[fontsize=\small,samepage=true,frame=leftline,framesep=5mm,framerule=1mm]
We got 82.33% correct on the test data.
\end{Verbatim}

Not too shabby.

